Una vez definido el examen, el coordinador convoca a los alumnos
mediante una \texttt{ExamConvocation} que enumera los identificadores
de cada alumno citado a cada modelo. Esta lista es la fuente de verdad
contra la que el servicio \texttt{matching} empareja la identificación
extraída por \acs{ocr} de cada página
(\cref{SS:DES-INGEST-ASSIGNMENT}).

La asignación de los correctores se modela mediante reglas declarativas
(\texttt{AssignmentRule}). Cada regla especifica un subconjunto de
instancias (por grupo, por modelo) y un subconjunto de problemas, y
los asigna a uno o varios correctores. Múltiples reglas pueden coexistir
sobre un mismo examen, lo que permite configuraciones flexibles como
``el problema 3 lo corrigen los profesores X e Y, el resto cualquier
corrector con permisos sobre la asignatura''. El servicio de asignación
combina todas las reglas activas y genera, para cada instancia, la
lista concreta de correctores autorizados a cada problema.
